\mxnewpage\section*{EV3-A -- Blanc 2}\label{EV3-A}
\addcontentsline{toc}{section}{EV3-A -- Blanc 2}
\textbf{Sujet réservé au professeur}. Programme étudié des chapitres 1 à 10.\par Durée : 2 h. Sans calculatrice. Note sur 20. 1–3 avril 2027.
\subsection*{Partie 1 -- Automatismes (6 points)}
Douze questions, 0{,}5 point chacune. Une seule réponse correcte ; aucune justification demandée. Une réponse incorrecte ou absente rapporte 0, sans retrait de points.
\begin{enumerate}
\item Calculer $3+2\times3$.\par A. $9$\quad B. $15$\quad C. $8$\quad D. $12$
\item Calculer $\frac12+\frac13$.\par A. $\frac{2}{5}$\quad B. $\frac{5}{6}$\quad C. $\frac{1}{6}$\quad D. $\frac{5}{3}$
\item Calculer $\dfrac{2^3}{2^{1}}$.\par A. $4$\quad B. $1$\quad C. $8$\quad D. $2$
\item Une valeur de 50 augmente de 20 \%. Valeur finale ?\par A. $10$\quad B. $40$\quad C. $60$\quad D. $70$
\item Une valeur de 80 est obtenue après une baisse de 20 \%. Valeur initiale ?\par A. $96$\quad B. $60$\quad C. $100$\quad D. $64$
\item Hausse de 25 \% puis baisse de 20 \% : coefficient global ?\par A. $\frac{4}{5}$\quad B. $\frac{21}{20}$\quad C. $\frac{9}{10}$\quad D. $1$
\item La solution de $2x+3=0$ vaut\par A. $\frac{3}{2}$\quad B. $\frac{-3}{2}$\quad C. $-6$\quad D. $6$
\item La pente entre $(0;3)$ et $(2;9)$ vaut\par A. $6$\quad B. $-3$\quad C. $2$\quad D. $3$
\item La moyenne de $3, 5, 7$ vaut\par A. $6$\quad B. $15$\quad C. $3$\quad D. $5$
\item Un secteur de $90^\circ$ représente quelle fraction du total ?\par A. $\frac{1}{3}$\quad B. $\frac{1}{4}$\quad C. $\frac{3}{4}$\quad D. $\frac{1}{2}$
\item Une probabilité vaut $\dfrac{2}{5}$. Son complément vaut\par A. $\frac{3}{5}$\quad B. $\frac{5}{2}$\quad C. $\frac{2}{5}$\quad D. $\frac{1}{5}$
\item Pour $f(x)=x^2+3$, $f(-2)$ vaut\par A. $12$\quad B. $7$\quad C. $5$\quad D. $-1$
\end{enumerate}
\mxnewpage\subsection*{Partie 2 -- Exercices indépendants (14 points)}
Toute démarche doit être justifiée. Les résultats exacts sont attendus sauf indication contraire.
\mxneed{12}
\subsection*{Exercice 1 -- Loterie et dispersion (5 points)}
Une loterie verse un gain brut $G$ de 0, 4 ou 12 unités monétaires avec les probabilités respectives $\dfrac12$, $\dfrac38$ et $\dfrac18$. Pour participer, on paie 3 unités. On note $X=G-3$ le gain net.
\begin{enumerate}
\item Donner la loi du gain net X. \hfill(1 pt)
\item Calculer espérance et variance de X. \hfill(2 pts)
\item On double les gains bruts et la mise. Que deviennent moyenne et écart type du gain net ? \hfill(2 pts)
\end{enumerate}
\mxneed{12}
\subsection*{Exercice 2 -- Dérivée et tangente horizontale (5 points)}
La fonction $f$ est définie sur $\R$ par $f(x)=(2x-3)e^x$. On étudie son extremum et la tangente à sa courbe au point d’abscisse 0. Toutes les valeurs sont attendues sous forme exacte.
\begin{enumerate}
\item Calculer $f(0)$, puis $f\prime(x)$ et $f\prime(0)$. \hfill(1 pt)
\item Étudier les variations de $f$ sur $\R$ et donner son extremum exact. \hfill(2 pts)
\item Déterminer l’équation de la tangente au point d’abscisse 0. Préciser si son point de contact est au-dessus ou au-dessous de l’axe des abscisses. \hfill(2 pts)
\end{enumerate}
\mxneed{12}
\subsection*{Exercice 3 -- Angle et longueur (4 points)}
Dans un repère orthonormé, on donne les points $A(0;0)$, $B(4;0)$ et $C(2;2\sqrt3)$. On cherche la nature du triangle $ABC$ en confrontant longueurs et produit scalaire.
\begin{enumerate}
\item Calculer AB et AC. \hfill(1 pt)
\item Déterminer le cosinus puis l’angle BAC en radians. \hfill(1 pt)
\item Déterminer BC et la nature du triangle. \hfill(2 pts)
\end{enumerate}
